\section*{Annexes}



\begin{table}[ht]
\centering
\footnotesize
\setlength{\tabcolsep}{4pt}
\begin{tabular}{llll}
\hline
\textbf{Satellite} & \textbf{Responsable} & \textbf{Période} & \textbf{Référence} \\
\hline
IMP 1 & Ness & 27/11/1963--15/02/1964 & Ness et al., 1964 \\
IMP 3 & Ness & 01/06/1965--29/01/1967 & Ness et al., 1964 \\
AIMP 1 & Ness & 04/07/1966--13/07/1968 & Behannon et al., 1968 \\
IMP 4 & Ness & 26/05/1967--27/12/1968 & Fairfield, 1969 \\
AIMP 2 & Ness & 26/07/1967--10/11/1969 & Ness et al., 1967 \\
HEOS 1 & Hedgecock & 11/12/1968--28/10/1975 & Hedgecock, 1975 \\
IMP 5 & Ness & 21/06/1969--26/10/1972 & Fairfield \& Ness, 1972 \\
IMP 6 & Ness & 14/03/1971--21/07/1974 & Fairfield, 1974 \\
IMP 7 & Ness & 26/09/1972--03/04/1973 & Mish \& Lepping, 1976 \\
IMP 8 & Ness & 30/10/1973--12/05/2000 & Mish \& Lepping, 1976 \\
ISEE 3 & E. Smith & 14/08/1978--21/12/1982 & Frandsen et al., 1978 \\
Prognoz 10 & Yeroshenko & 27/04/1985--04/11/1985 & Styazhkin et al., 1985 \\
Geotail & Nagai & 08/05/1995--29/09/1997 & -- \\
Wind & Lepping & 21/11/1994--en cours & Lepping et al., 1995 \\
ACE & Ness & 06/02/1998--en cours & Smith et al., 1998 \\
\hline
\end{tabular}
\caption{Satellites fournissant le champ magnétique interplanétaire (IMF) d'OMNI (\cite{SPDF_OMNIWeb}).}
\label{tab:omni_imf}
\end{table}

\begin{table}[ht]
\centering
\footnotesize
\setlength{\tabcolsep}{4pt}
\begin{tabular}{llll}
\hline
\textbf{Satellite} & \textbf{Responsable(s)} & \textbf{Période} & \textbf{Référence} \\
\hline
IMP 1 & Bridge & 27/11/1963--22/02/1964 & Bridge et al., 1965 \\
Vela (fusionné) & Bame & 21/07/1964--18/03/1971 & Bame et al., 1971 \\
Vela 3 & Bame & 26/07/1965--13/11/1967 & Hundhausen et al., 1967 \\
AIMP 1 & Bridge & 06/07/1966--23/09/1969 & Lyon et al., 1968 \\
IMP 4 & Ogilvie & 03/06/1967--16/12/1967 & Ogilvie et al., 1968 \\
AIMP 2 & Bridge & 28/07/1967--03/07/1968 & Lyon et al., 1967 \\
OGO 5 & Neugebauer & 05/03/1968--29/04/1971 & Neugebauer, 1970 \\
HEOS 1 & Bonetti & 11/12/1968--15/04/1970 & Bonetti et al., 1969 \\
IMP 6 & Bame & 18/03/1971--21/07/1974 & Feldman et al., 1973 \\
IMP 7 & Bame & 06/10/1972--29/09/1978 & Asbridge et al., 1976 \\
IMP 7 & Bridge & 03/01/1975--20/09/1978 & Lazarus et al., 1998 \\
IMP 8 & Bame & 04/11/1973--16/07/2000 & Asbridge et al., 1976 \\
IMP 8 & Bridge & 05/12/1973--26/07/2001 & Lazarus et al., 1998 \\
ISEE 1 & Bame & 30/10/1977--19/12/1979 & Bame et al., 1978b \\
ISEE 3 & Bame & 16/08/1978--12/10/1982 & Bame et al., 1978a \\
Geotail & L. Frank & 08/05/1995--07/12/2006 & -- \\
Wind & Lazarus et Kasper & 01/01/1995--en cours & Kasper, 2002 \\
ACE & McComas, R. Skoug & 05/02/1998--en cours & McComas et al., 1998 \\
\hline
\end{tabular}
\caption{Satellites fournissant le plasma (vent solaire) d'OMNI (\cite{SPDF_OMNIWeb}).}
\label{tab:omni_plasma}
\end{table}



% Nécessite le paquet longtable dans le préambule : \usepackage{longtable}

%==================================================================
% ANNEXE -- LRO (version horaire)
%==================================================================
\begingroup
\footnotesize
\setlength{\tabcolsep}{3pt}
\renewcommand{\arraystretch}{1.1}
\begin{longtable}{r p{5.0cm} p{2.0cm} l l p{3.3cm}}
\caption{Paramètres de la version horaire d'OMNI (LRO) et valeurs manquantes (\cite{OMNI_LRO_doc}).}
\label{tab:annexe_lro}\\
\hline
\textbf{N°} & \textbf{Paramètre} & \textbf{Unité} & \textbf{Format} & \textbf{Valeur manquante} & \textbf{Remarque} \\
\hline
\endfirsthead
\multicolumn{6}{l}{\textit{Suite du tableau~\ref{tab:annexe_lro}}} \\
\hline
\textbf{N°} & \textbf{Paramètre} & \textbf{Unité} & \textbf{Format} & \textbf{Valeur manquante} & \textbf{Remarque} \\
\hline
\endhead
\hline
\multicolumn{6}{r}{\textit{Suite page suivante}} \\
\endfoot
\hline
\endlastfoot
1 & Année & -- & I4 & -- & 1963, 1964, \dots \\
2 & Jour de l'année & jour & I4 & -- & 1er janvier = 1 \\
3 & Heure & h & I3 & -- & 0 à 23 ; « 1 » = moyenne de 01:00 à 02:00 \\
4 & Numéro de rotation de Bartels & -- & I5 & 9999 & \\
5 & Identifiant du satellite (champ magnétique) & -- & I3 & 99 & \\
6 & Identifiant du satellite (plasma) & -- & I3 & 99 & \\
7 & Nombre de points dans la moyenne du champ magnétique & -- & I4 & 999 & \\
8 & Nombre de points dans la moyenne du plasma & -- & I4 & 999 & \\
9 & Module moyen du champ $\langle F \rangle$ & nT & F6.1 & 999.9 & Moyenne scalaire \\
10 & Module du vecteur champ moyen $|\langle B \rangle|$ & nT & F6.1 & 999.9 & \\
11 & Angle de latitude du vecteur champ moyen & deg & F6.1 & 999.9 & Repère GSE \\
12 & Angle de longitude du vecteur champ moyen & deg & F6.1 & 999.9 & Repère GSE \\
13 & $B_x$ (GSE) & nT & F6.1 & 999.9 & \\
14 & $B_y$ (GSE) & nT & F6.1 & 999.9 & \\
15 & $B_z$ (GSE) & nT & F6.1 & 999.9 & \\
16 & $B_y$ (GSM) & nT & F6.1 & 999.9 & \\
17 & $B_z$ (GSM) & nT & F6.1 & 999.9 & \\
18 & $\sigma_{|B|}$ : écart-type du module (mot 10) & nT & F6.1 & 999.9 & Écart-type RMS \\
19 & $\sigma_B$ : écart-type du vecteur champ & nT & F6.1 & 999.9 & Écart-type RMS \\
20 & $\sigma_{B_x}$ (GSE) & nT & F6.1 & 999.9 & Écart-type RMS \\
21 & $\sigma_{B_y}$ (GSE) & nT & F6.1 & 999.9 & Écart-type RMS \\
22 & $\sigma_{B_z}$ (GSE) & nT & F6.1 & 999.9 & Écart-type RMS \\
23 & Température des protons & K & F9.0 & 9999999. & \\
24 & Densité de protons & cm$^{-3}$ & F6.1 & 999.9 & \\
25 & Vitesse du vent solaire & km/s & F6.0 & 9999. & Scalaire \\
26 & Longitude du flot $\phi_V$ & deg & F6.1 & 999.9 & \\
27 & Latitude du flot $\theta_V$ & deg & F6.1 & 999.9 & \\
28 & Rapport alpha/protons $N_a/N_p$ & -- & F6.3 & 9.999 & \\
29 & Pression dynamique & nPa & F6.2 & 99.99 & \\
30 & $\sigma_T$ & K & F9.0 & 9999999. & \\
31 & $\sigma_n$ & cm$^{-3}$ & F6.1 & 999.9 & \\
32 & $\sigma_V$ & km/s & F6.0 & 9999. & \\
33 & $\sigma_{\phi_V}$ & deg & F6.1 & 999.9 & \\
34 & $\sigma_{\theta_V}$ & deg & F6.1 & 999.9 & \\
35 & $\sigma_{N_a/N_p}$ & -- & F6.3 & 9.999 & \\
36 & Champ électrique & mV/m & F7.2 & 999.99 & \\
37 & Beta du plasma & -- & F7.2 & 999.99 & \\
38 & Nombre de Mach d'Alfvén & -- & F6.1 & 999.9 & \\
39 & Indice Kp $\times$ 10 & -- & I3 & 99 & Indice sur 3 h (GFZ) \\
40 & Nombre de taches solaires $R$ (version 2) & -- & I4 & 999 & Valeur journalière \\
41 & Indice Dst & nT & I6 & 99999 & Kyoto \\
42 & Indice AE & nT & I5 & 9999 & Kyoto \\
43 & Flux de protons $>1$~MeV & cm$^{-2}$ s$^{-1}$ sr$^{-1}$ & F10.2 & 999999.99 & \\
44 & Flux de protons $>2$~MeV & cm$^{-2}$ s$^{-1}$ sr$^{-1}$ & F9.2 & 99999.99 & \\
45 & Flux de protons $>4$~MeV & cm$^{-2}$ s$^{-1}$ sr$^{-1}$ & F9.2 & 99999.99 & \\
46 & Flux de protons $>10$~MeV & cm$^{-2}$ s$^{-1}$ sr$^{-1}$ & F9.2 & 99999.99 & \\
47 & Flux de protons $>30$~MeV & cm$^{-2}$ s$^{-1}$ sr$^{-1}$ & F9.2 & 99999.99 & \\
48 & Flux de protons $>60$~MeV & cm$^{-2}$ s$^{-1}$ sr$^{-1}$ & F9.2 & 99999.99 & \\
49 & Indicateur de contribution de la magnétosphère aux flux & -- & I3 & 0 & De 6 (aucune contribution) à 0 (tous les canaux) ; $-1$ : non vérifié \\
50 & Indice ap & nT & I4 & 999 & Indice sur 3 h (GFZ) \\
51 & Indice F10.7 & $10^{-22}$ W m$^{-2}$ Hz$^{-1}$ & F6.1 & 999.9 & Valeur journalière \\
52 & Indice PC(N) & -- & F6.1 & 999.9 & DTU, Copenhague \\
53 & Indice AL & nT & I6 & 99999 & Kyoto \\
54 & Indice AU & nT & I6 & 99999 & Kyoto \\
55 & Nombre de Mach magnétosonique & -- & F5.1 & 99.9 & \\
56 & Lyman-alpha solaire & W/m$^2$ & F9.6 & 0.999999 & Valeur journalière \\
57 & Quasi-invariant des protons QI & -- & F7.4 & 9.9999 & Énergie magnétique / énergie cinétique \\
\end{longtable}
\endgroup

\textit{Pour une valeur manquante, OMNI écrit un espace suivi de 9, qui remplit tout le format du mot. Par exemple, le format F6.1 donne 999.9 \cite{OMNI_LRO_doc}.}

%==================================================================
% ANNEXE -- HRO (1 min et 5 min)
%==================================================================
\begingroup
\footnotesize
\setlength{\tabcolsep}{3pt}
\renewcommand{\arraystretch}{1.1}
\begin{longtable}{r p{5.0cm} p{2.0cm} l l p{3.3cm}}
\caption{Paramètres des versions 1 min et 5 min d'OMNI (HRO) et valeurs manquantes (\cite{OMNI_HRO_doc}).}
\label{tab:annexe_hro}\\
\hline
\textbf{N°} & \textbf{Paramètre} & \textbf{Unité} & \textbf{Format} & \textbf{Valeur manquante} & \textbf{Remarque} \\
\hline
\endfirsthead
\multicolumn{6}{l}{\textit{Suite du tableau~\ref{tab:annexe_hro}}} \\
\hline
\textbf{N°} & \textbf{Paramètre} & \textbf{Unité} & \textbf{Format} & \textbf{Valeur manquante} & \textbf{Remarque} \\
\hline
\endhead
\hline
\multicolumn{6}{r}{\textit{Suite page suivante}} \\
\endfoot
\hline
\endlastfoot
1 & Année & -- & I4 & -- & \\
2 & Jour de l'année & jour & I4 & -- & 1 à 365 ou 366 \\
3 & Heure & h & I3 & -- & 0 à 23 \\
4 & Minute & min & I3 & -- & 0 à 59, début de la moyenne \\
5 & Identifiant du satellite (champ magnétique) & -- & I3 & 99 & ACE 71, Geotail 60, IMP~8 50, Wind 51 \\
6 & Identifiant du satellite (plasma) & -- & I3 & 99 & Mêmes identifiants \\
7 & Nombre de points dans la moyenne du champ magnétique & -- & I4 & 999 & \\
8 & Nombre de points dans la moyenne du plasma & -- & I4 & 999 & \\
9 & Pourcentage de points interpolés & \% & I4 & 999 & Normale au front de phase interpolée \\
10 & Décalage temporel & s & I7 & 999999 & \\
11 & Écart-type du décalage temporel & s & I7 & 999999 & RMS \\
12 & Écart-type de la normale au front de phase & -- & F6.2 & 99.99 & Absent en 5 min (valeur manquante) \\
13 & Temps entre observations (DBOT1) & s & I7 & 999999 & Une valeur négative indique une arrivée dans le désordre \\
14 & Module du champ magnétique & nT & F8.2 & 9999.99 & \\
15 & $B_x$ (GSE, GSM) & nT & F8.2 & 9999.99 & \\
16 & $B_y$ (GSE) & nT & F8.2 & 9999.99 & \\
17 & $B_z$ (GSE) & nT & F8.2 & 9999.99 & \\
18 & $B_y$ (GSM) & nT & F8.2 & 9999.99 & Calculé après le décalage \\
19 & $B_z$ (GSM) & nT & F8.2 & 9999.99 & Calculé après le décalage \\
20 & Écart-type du module du champ & nT & F8.2 & 9999.99 & RMS \\
21 & Écart-type du vecteur champ & nT & F8.2 & 9999.99 & RMS \\
22 & Vitesse du vent solaire & km/s & F8.1 & 99999.9 & \\
23 & $V_x$ (GSE) & km/s & F8.1 & 99999.9 & \\
24 & $V_y$ (GSE) & km/s & F8.1 & 99999.9 & \\
25 & $V_z$ (GSE) & km/s & F8.1 & 99999.9 & \\
26 & Densité de protons & cm$^{-3}$ & F7.2 & 999.99 & \\
27 & Température & K & F9.0 & 9999999. & \\
28 & Pression dynamique & nPa & F6.2 & 99.99 & \\
29 & Champ électrique & mV/m & F7.2 & 999.99 & \\
30 & Beta du plasma & -- & F7.2 & 999.99 & \\
31 & Nombre de Mach d'Alfvén & -- & F6.1 & 999.9 & \\
32 & Position du satellite, $X$ (GSE) & R$_E$ & F8.2 & 9999.99 & \\
33 & Position du satellite, $Y$ (GSE) & R$_E$ & F8.2 & 9999.99 & \\
34 & Position du satellite, $Z$ (GSE) & R$_E$ & F8.2 & 9999.99 & \\
35 & Nez du bow shock, $X$ (GSE) & R$_E$ & F8.2 & 9999.99 & \\
36 & Nez du bow shock, $Y$ (GSE) & R$_E$ & F8.2 & 9999.99 & \\
37 & Nez du bow shock, $Z$ (GSE) & R$_E$ & F8.2 & 9999.99 & \\
38 & Indice AE & nT & I6 & 99999 & Kyoto \\
39 & Indice AL & nT & I6 & 99999 & Kyoto \\
40 & Indice AU & nT & I6 & 99999 & Kyoto \\
41 & Indice SYM/D & nT & I6 & 99999 & Kyoto \\
42 & Indice SYM/H & nT & I6 & 99999 & Kyoto \\
43 & Indice ASY/D & nT & I6 & 99999 & Kyoto \\
44 & Indice ASY/H & nT & I6 & 99999 & Kyoto \\
45 & Indice PC(N) & -- & F7.2 & 999.99 & Université technique du Danemark. Dans la version modifiée : rapport $N_a/N_p$ (F7.3, 9.999) \\
46 & Nombre de Mach magnétosonique & -- & F5.1 & 99.9 & \\
47 & Flux de protons $>10$~MeV$^{*}$ & cm$^{-2}$ s$^{-1}$ sr$^{-1}$ & F9.2 & 99999.99 & 5 min seulement \\
48 & Flux de protons $>30$~MeV$^{*}$ & cm$^{-2}$ s$^{-1}$ sr$^{-1}$ & F9.2 & 99999.99 & 5 min seulement \\
49 & Flux de protons $>60$~MeV$^{*}$ & cm$^{-2}$ s$^{-1}$ sr$^{-1}$ & F9.2 & 99999.99 & 5 min seulement \\
\end{longtable}
\endgroup

\textit{Les mots 47 à 49 (flux GOES) n'existent que dans la version 5 minutes. $^{*}$ L'unité n'est pas précisée dans la documentation HRO : elle est reprise de la documentation LRO. Comme pour LRO, une valeur manquante est écrite avec des 9 qui remplissent le format \cite{OMNI_HRO_doc}.}